\documentclass[10pt,a4paper]{amsart}
\usepackage[margin=19mm]{geometry}
\usepackage{amsmath,amssymb,mathtools}
\usepackage[scr=rsfs,cal=euler]{mathalfa}
\usepackage{xfrac}
\usepackage[unicode,hidelinks,bookmarks=false]{hyperref}
\newcommand{\ie}{i.e.\ }
\newcommand{\cf}{cf.\ }
\newcommand{\resp}{respectively\ }
\newcommand{\Subsection}{\subsection}
\providecommand{\nobreakdash}{\nobreak\hbox{-}}
\setlength{\emergencystretch}{2em}
\multlinegap0pt
\usepackage{bm}
\usepackage{bbm}
\usepackage{dsfont}
\usepackage{xspace}
\usepackage{cancel}
\usepackage{mathtools}
\usepackage{amsthm}
\makeatletter
\@ifundefined{theorem}{\newtheorem{theorem}{Theorem}}{}
\@ifundefined{corollary}{\newtheorem{corollary}[theorem]{Corollary}}{}
\@ifundefined{lemma}{\newtheorem{lemma}[theorem]{Lemma}}{}
\@ifundefined{proposition}{\newtheorem{proposition}[theorem]{Proposition}}{}
\makeatother
\title[The Prescribed-Vertex Semidegree Threshold for Directed $3q$]{The Prescribed-Vertex Semidegree Threshold for Directed $3q$-Cycles in Oriented Graphs}
\author{Zhenhua Lyu}
\date{Source: arXiv:2608.20048v1 (CC BY 4.0)}
\begin{document}
\maketitle

\noindent\emph{Source and licence.} The Prescribed-Vertex Semidegree Threshold for Directed $3q$-Cycles in Oriented Graphs. Version arXiv:2608.20048v1 (CC BY 4.0), distributed under the Creative Commons Attribution 4.0 International licence, as recorded in the arXiv licence metadata for this exact version and shown on the abstract page https://arxiv.org/abs/2608.20048v1. Full rights notice: licenses/arxiv-2608.20048-CC-BY-4.0.txt.\par\smallskip

\noindent\textit{Editorial scope.} Counted unit: the Proposition whose source label is \texttt{prop:first-port} in the article (source0.tex lines 1090--1093, source label \texttt{prop:first-port}), delivered together with the article's own complete original proof of it (source0.tex lines 1095--1233, one \texttt{proof} environment of 139 lines). No mathematics is written, repaired or re-derived: statement and proof are the article's own text line for line. Required context retained uncounted: lem:extension-deletion (lines 526--560); eq:bookkeeping (lines 751--754); cor:stability (lines 996--1006); eq:row-meets-C (lines 1019--1022). Every definition, display and label of those lines is retained unchanged. Omitted from this package and not redistributed: 1--525, 561--750, 755--995, 1007--1018, 1023--1089, 1094--1094, 1234--1642. Nothing inside a retained excerpt is omitted or altered: every retained body is delivered exactly as the article's lines are written, so no omission receipt of any kind accompanies this package. Citations: the retained excerpts carry no citation command at all, so no bibliography entry of the article is transcribed and no reference is redistributed. Reference adaptation: none was needed and none was made; every reference inside the delivered unit resolves inside it, so no unresolved reference or citation is produced. Printed slips kept literal: no printed word, symbol, hypothesis or number of the counted statement or of its proof was altered. Layout and class: the article's own class is \texttt{amsart} with packages including fontenc, lmodern, amsmath, amssymb, mathtools, enumitem, microtype, needspace, hyperref; the wrapper is the lane's pinned \texttt{amsart} editorial wrapper at 10pt on A4, which restates the theorem-like environments the retained text needs and adds the article's own macro definitions that the retained text needs (none), with extra wrapper packages (bm, bbm, dsfont, xspace, cancel, mathtools). The delivered numbering is the unit's own. Assets: no retained span contains a figure environment, a graphics-inclusion command, a PDF-page inclusion, a table or a TikZ picture; no image file is copied into this package, the package declares no asset dependency and no image file is redistributed with it. Licence: version arXiv:2608.20048v1 of this article is distributed under the Creative Commons Attribution 4.0 International licence, as recorded in the arXiv licence metadata for this exact version and shown on the abstract page https://arxiv.org/abs/2608.20048v1. A full rights notice is delivered at licenses/arxiv-2608.20048-CC-BY-4.0.txt. Characters: every retained excerpt is pure ASCII, so the delivered engine, XeLaTeX with its default Latin Modern text font, reports no missing character.
\begin{lemma}\label{lem:extension-deletion}
Let $G$ be an oriented graph.
\begin{enumerate}
\item Let $d$ and $L$ be nonnegative integers.  If
$\delta^+(G)\ge d$, $F\subseteq V(G)$, $v\notin F$, and
\[
  |F|+L\le d,
\]
then $G$ contains a simple path of length $L$ starting at $v$ and otherwise avoiding $F$.
\item Let $G$ have order $n$ and minimum semidegree at least
$\lceil n/3\rceil$, and let $R\ge0$ be an integer.  If
\[
  n\ge15R+7
\]
and $H$ is obtained by deleting at most $R$ vertices, then every
ordered pair of distinct vertices of $H$ is joined in $H$ by a path
of length $3$, $4$, or $5$.
\item Let $G$ have order $n$ and minimum semidegree at least
$\lceil n/3\rceil$.  Let $\ell$ and $B$ be integers with
$\ell\ge B\ge6$, let $x,a$ be distinct
vertices, and let $Z\subseteq V(G)\setminus\{x,a\}$.  Suppose that,
for each $j\in\{B-5,B-4,B-3\}$, there is an $x$--$a$ path $Q_j$ of
length $j$ whose internal vertices lie in $Z$.  Put
\[
  R=|Z|+\ell-B.
\]
If
\[
  |Z|+1+\ell-B\le\left\lceil\frac n3\right\rceil
  \qquad\text{and}\qquad
  n\ge15R+7,
\]
then $x$ lies on a copy of $C_\ell$.
\end{enumerate}
\end{lemma}

\begin{equation}\label{eq:bookkeeping}
  d_B^+(b)\ge s_b,\qquad d_B^-(b)\ge t_b,\qquad
  \sum_{b\in B}s_b=\sum_{b\in B}t_b=m^2.
\end{equation}

\begin{corollary}\label{cor:stability}
If a row family indexed by $I$, $|I|\ge4$, has no terminal-safe three-chain, then either
\begin{enumerate}
\item one row type has multiplicity at least $|I|-1$, or
\item
\[
  \left|\left(\bigcup_{i\in I}O_i\right)
  \setminus\left(\bigcap_{i\in I}O_i\right)\right|\le3.
\]
\end{enumerate}
\end{corollary}

\begin{equation}\label{eq:row-meets-C}
  O(a)\cap C\ne\varnothing
  \qquad\text{for every }a\in A.
\end{equation}

\begin{proposition}\label{prop:first-port}
Let $q\ge3$, put $\ell=3q$, and suppose $n=3m\ge n_0(q)$.
Then $x$ lies on a copy of $C_\ell$.
\end{proposition}

\begin{proof}
Since $n=3m\ge n_0(q)$,
\[
  m\ge
  \begin{cases}
    33,&q=3,\\
    15q-2,&q\ge4.
  \end{cases}
\]
In particular,
\begin{equation}\label{eq:first-port-numerics}
  m\ge\ell,\qquad m\ge2q+11,\qquad m-4\ge q.
\end{equation}
Let $W$ be the nonextendable-endpoint set and put $A_0=A\setminus W$.
Thus $|W|\le3$ and every $a\in A_0$ admits a $D_3$-transition.  Apply
Corollary~\ref{cor:stability} to the rows indexed by $A_0$, with terminal
set $C$.

The proof follows the three outcomes in Corollary~\ref{cor:stability}.
A terminal-safe three-chain gives three possible initial lengths and is
closed by short linking.  In each of the two remaining outcomes,
\eqref{eq:bookkeeping} gives a dense bipartite graph, and K\"onig's
theorem supplies the required matching.

\smallskip
\noindent\emph{Case 1: a terminal-safe three-chain.}
Suppose first that there is a terminal-safe three-chain.  It yields distinct vertices with
\begin{equation}\label{eq:safe-chain}
  x\to a_0\to b_0\to a_1\to b_1\to a_2\to b_2\to a_3\to d\to x.
\end{equation}
If $q=3$, this is already a $C_9$.  Assume $q\ge4$.  The suffixes of~\eqref{eq:safe-chain} give $x$--$a_3$ paths
\[
  Q_3=xa_2b_2a_3,\qquad
  Q_5=xa_1b_1a_2b_2a_3
\]
of lengths three and five.  Since $a_3\notin W$, take a transition
\[
  a'\to u\to v\to a_3.
\]
Then
\[
  Q_4=xa'uva_3
\]
is a simple path of length four.  Different candidate paths may intersect; only one will eventually be used.  Let $Z$ be the union of their internal vertices.  Then $|Z|\le7$.

Since $|Z|\le7$, we have
\[
  |Z|+1+\ell-8\le\ell\le m,
  \qquad
  15(|Z|+\ell-8)+7\le15(\ell-1)+7=n_0(q).
\]
Lemma~\ref{lem:extension-deletion}(iii), with $B=8$, now closes one of
$Q_3,Q_4,Q_5$ to a $C_\ell$.

\smallskip
\noindent\emph{Case 2: a dominant row.}
Suppose that a row type $S$ occurs at least
\[
  |A_0|-1\ge m-4
\]
times.  For every $b\in S$, we have $s_b\ge m-4$, so
\begin{equation}\label{eq:heavy-row}
  e(S,B\setminus S)
  \ge m(m-4)-\binom m2
  =\frac{m^2-7m}{2}.
\end{equation}
Choose $d\in S\cap C$, which exists by~\eqref{eq:row-meets-C}, and
form the bipartite graph with parts
\[
  S\setminus\{d\}
  \quad\text{and}\quad
  (B\setminus S)\setminus\{x\},
\]
joining $b$ to $c$ precisely when $b\to c$ in $G$.  Deleting the
source $d$ and the target $x$ removes at most $2m$ arcs
from~\eqref{eq:heavy-row}, so this bipartite graph has at least
$(m^2-11m)/2$ edges.  If its maximum matching had size at most $q-2$,
K\"onig's theorem would give a vertex cover of size at most $q-2$.
Every vertex of the bipartite graph has degree at most $m$, so the cover
would meet at most $(q-2)m$ edges.  By~\eqref{eq:first-port-numerics},
\[
  \frac{m^2-11m}{2}-(q-2)m
  =\frac m2(m-2q-7)>0,
\]
a contradiction.  Take a matching
\[
  b_i\to c_i\qquad(0\le i\le q-2)
\]
and distinct $S$-type vertices $a_0,\ldots,a_{q-1}$.  These choices are
available because $m-4\ge q$ by~\eqref{eq:first-port-numerics}.  Then
\[
  x\to a_0\to b_0\to c_0\to a_1\to\cdots
  \to b_{q-2}\to c_{q-2}\to a_{q-1}\to d\to x
\]
is a $C_{3q}$.

\smallskip
\noindent\emph{Case 3: variation on at most three points.}
Suppose the active variation has size at most three.  Put
\[
  K=\bigcap_{a\in A_0}O(a),\qquad
  V=\left(\bigcup_{a\in A_0}O(a)\right)\setminus K,
  \qquad L=B\setminus(K\cup V).
\]
Write $v=|V|$ and $|K|=m-r$.  Every row is $K\cup P_a$ with $P_a\subseteq V$ and $|P_a|=r$.  If $v>0$, each point of $V$ appears in some but not all rows, so $1\le r\le v-1$.  Hence
\[
  (v,r)\in\{(0,0),(2,1),(3,1),(3,2)\}.
\]
Every $b\in K$ belongs to all rows indexed by $A_0$, and therefore $s_b\ge m-3$.  Since $|L|=m+r-v\le m$, we obtain
\begin{align}
e(K,L)
&\ge |K|(m-3)-\binom{|K|}{2}-|K|v\notag\\
&=\frac{(m-r)(m+r-5-2v)}2\notag\\
&\ge\frac{m^2-11m+10}{2}.
\label{eq:hub-count}
\end{align}
Fix $a_*\in A_0$ and $d\in O(a_*)\cap C$.  Since $x$ belongs to no
row, $x\in L$.  Form the bipartite graph whose left part is
$K\setminus\{d\}$ if $d\in K$ and is $K$ otherwise, whose right part
is $L\setminus\{x\}$, and whose edges are the arcs directed from left
to right.  Deleting the target $x$ and, when necessary, the source $d$
removes at most $2m$ arcs from~\eqref{eq:hub-count}; hence at least
$(m^2-15m+10)/2$ edges remain.  If there were no matching of size
$q-1$, K\"onig's theorem would give a vertex cover of size at most
$q-2$.  Both parts have size at most $m$, so such a cover would meet at
most $(q-2)m$ edges.  By~\eqref{eq:first-port-numerics},
\[
  m^2-(2q+11)m+10>0
\]
and hence there is a matching $b_i\to c_i$ of size $q-1$ from $K$ to
$L$, avoiding $d$ and $x$.  Since $|A_0|\ge m-3\ge q$, choose
distinct $a_0,\ldots,a_{q-1}\in A_0$ with $a_{q-1}=a_*$.  Since every
row contains $K$ and avoids $L$,
\[
  x\to a_0\to b_0\to c_0\to a_1\to\cdots
  \to b_{q-2}\to c_{q-2}\to a_{q-1}\to d\to x
\]
is again a $C_{3q}$.
\end{proof}

\end{document}
